\section{ضرورت انجام پژوهش}
% Reconstructed and edited from physical pages 26--27 of the source PDF.
خودکشی یک مسئلۀ مهم سلامت عمومی با پیامدهای فردی و اجتماعی گسترده است. سازمان جهانی بهداشت در گزارش مورد استناد، شمار سالانۀ جان‌باختگان را در حدود ۷۰۰ هزار نفر توصیف می‌کند و بر پیشگیری و دسترسی به حمایت تأکید دارد \cite{ref1}. این پدیده نتیجۀ تعامل عوامل ژنتیکی، روان‌شناختی، اجتماعی و محیطی است؛ ازاین‌رو توضیح یا پیش‌بینی آن با یک نشانۀ منفرد کافی نیست \cite{ref3}. شناسایی زودهنگام می‌تواند فرصت مداخله ایجاد کند، اما نیازمند داده و ابزار ارزیابی معتبر است.

اهمیت مسئله در گروه‌های سنی جوان نیز در آمارهای سلامت عمومی دیده می‌شود \cite{ref2,ref5}. شکل‌گیری هویت، فشارهای تحصیلی و شغلی، شروع زندگی مستقل و محدودیت مهارت‌های مقابله‌ای می‌توانند بخشی از زمینۀ آسیب‌پذیری باشند. در کنار آن، ترس از قضاوت و انگ بیماری روانی ممکن است مراجعه به خدمات رسمی را دشوار کند. برخی افراد تجربه‌ها و احساسات خود را در محیط‌های آنلاین بیان می‌کنند؛ این رفتار امکان مطالعۀ نشانه‌های زبانی را فراهم می‌کند، اما مجوز تفسیر قطعی وضعیت بالینی هر نویسنده نیست \cite{ref6,ref7}.

شبکه‌های اجتماعی و یادداشت‌های متنی، داده‌هایی با حجم زیاد و ساختار ناهمگون ایجاد می‌کنند. پژوهش‌های زبان‌شناختی به الگوهایی مانند استفاده از ضمیر اول‌شخص، واژگان مرتبط با تنهایی، ناامیدی یا مرگ و تغییر لحن توجه کرده‌اند \cite{ref8,ref9}. روش‌های پردازش زبان طبیعی و یادگیری ماشین می‌توانند این الگوها را در چارچوب یک مجموعه‌دادۀ برچسب‌گذاری‌شده تحلیل کنند. مزیت آن‌ها مقیاس‌پذیری تحلیل است؛ محدودیت آن‌ها نیز وابستگی به تعریف برچسب، بافت، زبان و دادۀ آموزشی است \cite{ref10}.

مجموعه‌دادۀ استاندارد و نمایندۀ همۀ کاربران به‌سادگی در دسترس نیست. متن‌های شبکۀ اجتماعی با مصاحبۀ بالینی تفاوت دارند و نمونه‌گیری از یک بستر خاص می‌تواند سوگیری ایجاد کند. تفاوت فرهنگ، واژگان، سن، زمان و دسترسی کاربران به خدمات نیز بر تعمیم مدل اثر دارد. بنابراین ارزیابی باید همراه با تعریف روشنِ داده، کنترل تکرار، تفکیک کاربران یا متن‌های هم‌پوشان و بیان حدود کاربرد انجام شود. یک عدد دقت بالا، این مسائل را خودکار حل نمی‌کند.

مانع دیگر، حساسیت شدید اطلاعات سلامت روان است \cite{ref4}. افشای متن می‌تواند با انگ اجتماعی، تبعیض یا پیامدهای شخصی همراه باشد. قواعدی مانند \en{GDPR} و \en{HIPAA} نیز اهمیت تعیین هدف، دسترسی مجاز، نگهداری مسئولانه و رضایت را نشان می‌دهند. نیاز پژوهشگر به دادۀ متنوع باید با مسئولیت حفاظت از داده و محدودیت استفادۀ ثانویه سازگار شود. این نیاز، انگیزه‌ای برای بررسی روش‌هایی است که ضرورت جابه‌جایی دادۀ خام را کاهش می‌دهند.

یادگیری مشارکتی یکی از این مسیرهاست. در شکل‌های رایج آن، آموزش روی دستگاه یا سازمان محلی انجام و اطلاعاتی مانند به‌روزرسانی مدل یا خروجی پیش‌بینی به یک هماهنگ‌کننده فرستاده می‌شود. نوع اطلاعات اشتراکی به طراحی الگوریتم وابسته است. در روش این پایان‌نامه، اشتراک در سطح احتمال پیش‌بینی صورت می‌گیرد و دادۀ آموزش کلاینت‌ها برای آموزش یک مدل واحد به سرور منتقل نمی‌شود. این طراحی با میانگین‌گیری پارامترهای مدل یکسان نیست.

باقی ماندن دادۀ آموزش در محل، تنها یکی از لایه‌های حفاظت است. احتمال خروجی و به‌روزرسانی نیز می‌توانند اطلاعاتی حمل کنند؛ مهاجم فعال می‌تواند هویت جعل کند، وزن را تغییر دهد یا یک پیام معتبر را دوباره ارسال کند. در نتیجه، محرمانگی، اصالت، تمامیت و کنترل بازپخش باید به‌صورت جداگانه بررسی شوند. رمزگذاری همریخت، ابزار انجام محاسبه در دامنۀ رمز و امضای دیجیتال، ابزار بررسی اصالت محموله است؛ هیچ‌کدام جای تمام کنترل‌های سامانه را نمی‌گیرد.

ضرورت پژوهش حاضر از کنار هم قرار گرفتن این نیازها ناشی می‌شود: ارزیابیِ بدون نشت برای تحلیل متن، تعریف دقیق توزیع دادۀ کلاینت‌ها و سنجش یک روش احراز اصالت سازگار با تجمیع جمعی. هدف، گزارش یک چارچوب پژوهشیِ روشن و قابل ممیزی است. ادعای کاربرد بالینی، محرمانگی مطلق در برابر دارندۀ کلید یا اجرای اثبات دانش صفر، از آزمایش‌های انجام‌شده نتیجه گرفته نمی‌شود.
